\maketitle
\tableofcontents

\chapter*{Dok.~311 -- Vier auf drei}
\addcontentsline{toc}{chapter}{Dok.~311 -- Vier auf drei}

\section*{Worum es geht}

Ein Gitter in $\mathbb{R}^4$ hat vier gleichberechtigte Richtungen.
Gemessen wird in drei. Wie kommt man von den vieren auf die dreie?

Es gibt genau drei Wege, und sie kosten sehr Verschiedenes. Diese Notiz
geht sie durch, verdichtet sie am Ende auf zwei Bedingungen und legt die
Bedingungen an drei Fälle an: die eigene Konstruktion, generische
Projektionskonstruktionen und die Stringtheorie.

\medskip
\noindent\textbf{Status:} Arbeitsnotiz. Nicht Teil der A-Serie.
Die Abschnitte zu den Wegen und zum Zerfallsbefund sind gerechnet
(\texttt{ffgft\_311\_vier\_auf\_drei.py}); das Kapitel
\enquote{Anwendung} stellt eine Frage und trifft keine Feststellung.

\chapter{Die drei Wege}

\section{Weg 1 --- Kompaktifizieren}

Die vierte Richtung wird ein Kreis vom Radius $R$. Kaluza 1921,
Klein 1926.

\textbf{Was es bringt.} Auf einem Kreis sind nur ganzzahlige Wicklungen
zulässig, und daraus entsteht ein Modenturm mit $m_n \propto n/R$. Der
Modenindex \emph{ist} eine Skala. Nichts muss angeheftet werden.

\textbf{Was es kostet.} $R$ ist ein freier Modul. Wer nicht sagt,
wodurch $R$ festliegt, hat eine Skala gewonnen und einen Parameter dazu
--- also nichts. Das ist der Preis, den die Stringtheorie mit ihrer
Landschaft zahlt.

Der Modul verschwindet nur, wenn $R$ an eine andere Größe gebunden
wird. In FFGFT geschieht das über $\tilde T \cdot m = 1$: sobald die
Masse feststeht, steht der Zyklus fest, und es bleibt nichts zu eichen.

\section{Weg 2 --- Projizieren}

Das 4D-Gitter wird auf einen 3D-Schnitt projiziert. Der Standardweg für
Quasikristalle, cut-and-project.

\textbf{Was es kostet.} Die Projektion hat einen Fensterparameter, und
die Gitterkennzahlen überleben sie nicht. Für D4 auf die ersten drei
Koordinaten:
\[
  24 \text{ kürzeste Vektoren} \;\longrightarrow\; 18 \text{ verschiedene Bilder.}
\]
Sechs Paare fallen zusammen. Was projiziert ankommt, ist kein Gitter
mehr mit Kusszahl~24, sondern ein aperiodisches Punktmuster mit anderen
Kennzahlen.

\section{Weg 3 --- Die vierte ist kein Raum}

Die vierte Koordinate trägt eine andere Größe: Zeit, Masse, Ladung,
einen inneren Index. Dann sind nur drei Richtungen Raum, und es ist
nichts zu reduzieren.

\textbf{Was es kostet.} Man muss sagen, \emph{was} sie trägt, und die
Aussage muss tragfähig sein. Das ist keine Formalie: eine vierte
Richtung, die \enquote{Masse} heißt, ohne dass etwas sie an eine Masse
bindet, ist nur ein Etikett.

\chapter{Weg 1 und Weg 3: einzeln beide unvollständig}

Die beiden sehen sich ähnlich, weil beide die vierte Richtung aus dem
Raum herausnehmen. Sie tun aber Verschiedenes, und der Unterschied ist
genau die Stelle, an der eine Konstruktion trägt oder nicht.

\textbf{Weg 1 macht die Richtung kompakt und lässt offen, was sie
trägt.} Reines Kaluza-Klein. Ein Kreis vom Radius $R$ erzwingt
ganzzahlige Wicklungen, daraus entsteht ein Modenturm
$m_n \propto n/R$, und der Modenindex ist eine Skala. Aber $R$ selbst
ist frei. Man hat eine Skala gewonnen und einen Parameter dazugekauft.

\textbf{Weg 3 sagt, was die Richtung trägt, und lässt offen, ob sie
kompakt ist.} Sie heißt \enquote{Masse} oder \enquote{Ladung} oder
\enquote{innerer Index}. Ohne Kompaktheit ist das ein Etikett: nichts
erzwingt ganzzahlige Werte, ohne Ganzzahligkeit gibt es keinen
Modenturm, und ohne Modenturm keine Skala.

\textbf{Beide sind für sich unvollständig, auf entgegengesetzte
Weise:}
\begin{center}
\begin{tabular}{ll}
Weg 1: & Mechanismus ohne Bedeutung \\
Weg 3: & Bedeutung ohne Mechanismus \\
\end{tabular}
\end{center}

\section{Was die Lücke schließt}

Beides zugleich, und die Bindung dazwischen.

\begin{center}
\begin{tabular}{lllp{4.6cm}}
\toprule
 & kompakt & gebunden & Ergebnis \\
\midrule
Kaluza-Klein & ja & nein & Skala vorhanden, aber freier Modul \\
bloßes Etikett & nein & dem Namen nach & kein Modenturm, keine Skala \\
\textbf{FFGFT} & \textbf{ja} & \textbf{ja, über $\tilde T \cdot m = 1$}
  & \textbf{Skala ohne Modul} \\
\bottomrule
\end{tabular}
\end{center}

Die Reziprozität schließt beide Lücken in einem Schritt. Der Modenturm
bekommt eine Bedeutung, weil die Richtung die Massenrichtung ist. Und
der freie Radius verschwindet, weil $R$ nicht mehr wählbar ist, sobald
$m$ feststeht --- der Zyklus \emph{ist} die Masse, nicht ein Behälter
für sie.

\section{Die schärfere Frage}

Damit ist die eigentliche Unterscheidung nicht \enquote{welcher der
drei Wege}, sondern:

\begin{quote}
Ist die vierte Richtung kompakt \textbf{und} an etwas gebunden --- oder
nur eines von beidem?
\end{quote}

Wer nur kompaktifiziert, hat einen Modul zu stabilisieren. Wer nur
benennt, hat ein Wort. Beides zusammen und aneinander gebunden ist der
einzige Fall, in dem eine Skala entsteht, ohne dass ein Parameter
nachkommt.

Das ist auch der Prüfstein für fremde Konstruktionen, und er ist in
einer Frage zu stellen: \emph{was legt den Radius fest?} Kommt darauf
keine Antwort, ist die Skala nur verschoben.

\chapter{Weg 2 genauer: was Projizieren voraussetzt und kostet}

Weg 2 unterscheidet sich von den anderen beiden nicht dem Grad nach,
sondern der Art nach.

\section{Projektion braucht einen Standpunkt außerhalb}

Kompaktifizieren ist eine Eigenschaft des Objekts: der Kreis ist da
oder er ist nicht da, unabhängig davon, wer hinsieht. Die vierte
Richtung als Massenrichtung zu deuten ist eine Aussage über die
Struktur.

Projizieren ist eine Abbildung \textbf{von außen}. Es braucht eine
Schnittrichtung und ein Fenster, und beide wählt der Beobachter. Sie
sagen nichts über das Gitter aus, sondern über den Zugriff darauf.

Das ist ein anderer Typ von freiem Parameter als der Radius bei Weg~1.
Ein Radius ist ein Merkmal des Objekts und lässt sich an eine andere
Größe des Objekts binden --- genau das leistet
$\tilde T \cdot m = 1$. Eine Schnittrichtung lässt sich nicht binden.
Sie lässt sich nur festlegen, und die Festlegung bleibt Setzung.

\section{Was mit den 24 Vektoren geschieht}

Bei Projektion auf die ersten drei Koordinaten:
\begin{center}
\begin{tabular}{lcl}
24 Vektoren & $\longrightarrow$ & 18 Bilder \\[4pt]
12 FCC-Vektoren ($x_4 = 0$) & $\longrightarrow$
  & 12 Bilder, Länge $\sqrt2$, unberührt \\
12 Wicklungsvektoren & $\longrightarrow$
  & \phantom{1}6 Bilder, Länge 1, paarweise verschmolzen \\
\end{tabular}
\end{center}

Denn $(1,0,0,+1)$ und $(1,0,0,-1)$ haben dasselbe Bild $(1,0,0)$. Der
Unterschied zwischen zwei entgegengesetzten Wicklungen verschwindet.

Und im Bild gibt es zwei verschiedene Längen, $1$ und $\sqrt2$. Das
Bild ist kein Gitter mit einheitlicher Kusszahl mehr; Packungsdichte im
üblichen Sinn verliert dort ihren Sinn.

\textbf{Das ist der eigentliche Verlust.} Beim Kompaktifizieren bleiben
die zwölf Wicklungsnachbarn unterscheidbar --- $+1$ und $-1$ sind
verschiedene Wicklungen und tragen verschiedene Quantenzahlen. Beim
Projizieren werden sie ununterscheidbar gemacht. Was aus dem Raum
herausführt, wird gelöscht statt umgedeutet.

\section{Wo die irrationalen Zahlen herkommen}

Bei rationaler Schnittrichtung, etwa einer Koordinatenachse, bleiben
endlich viele Bildlängen übrig --- hier zwei. Die Verkürzung der
Wicklungsvektoren um $1/\sqrt2$ ist zwar irrational, aber sie
wiederholt sich.

Beim eigentlichen Cut-and-Project ist die Schnittrichtung irrational
gewählt, und dann sind die Projektionslängen generisch irrational. Für
die Richtung $(1, \varphi, 0)$:
\begin{center}
\begin{tabular}{lcr}
$(1, 0, 0)$ & $\longrightarrow$ & $0{,}525731$ \\
$(0, 1, 0)$ & $\longrightarrow$ & $0{,}850651$ \\
$(1, 1, 0)$ & $\longrightarrow$ & $1{,}376382$ \\
$(1,-1, 0)$ & $\longrightarrow$ & $-0{,}324920$ \\
\end{tabular}
\end{center}

Jede Länge eine neue Zahl, und das Bild ist \textbf{dicht}. Erst das
Fenster macht es wieder diskret --- der zweite freie Parameter.

Genau so entstehen Quasikristalle. Und genau so kommt $\varphi$ dort
hinein:

\begin{quote}
Die Streckungsfaktoren sind irrational, ja. Aber sie sind Eigenschaften
der \textbf{gewählten Richtung}, nicht des Gitters. Wählt man anders,
sind es andere Zahlen.
\end{quote}

\section{Die Folge für fremde Konstruktionen}

Wer $\varphi$ über eine Cut-and-Project-Konstruktion hereinholt, hat es
hineingelegt. Es kommt nicht aus dem Gitter, sondern aus der Wahl der
Projektionsrichtung. Das ist zulässig, sofern es gesagt wird --- und es
ist etwas grundsätzlich anderes als eine Zahl, die aus der Struktur
folgt.

Der Prüfstein aus dem vorigen Kapitel bekommt damit eine zweite Frage
neben \emph{was legt den Radius fest?}:

\begin{quote}
\textbf{Wer hat die Schnittrichtung gewählt, und woraus?}
\end{quote}

Kommt darauf keine Antwort, sind alle Zahlen, die aus der Projektion
stammen, Eigenschaften des Beobachters.

\chapter{Die Umkehrung: wo \texorpdfstring{$\varphi$}{phi} gefunden
wird, ist die Richtung nicht mehr frei}

Wenn bei cut-and-project $\varphi$ erscheint, weil die Schnittrichtung
so gewählt wurde --- was gilt dann, wenn $\varphi$ nicht gewählt,
sondern \emph{gefunden} wird?

\section{Wo \texorpdfstring{$\varphi$}{phi} im Korpus steht}

A120 leitet $\theta = 2/9$ aus der Umverteilung der Lepton-Moden unter
einer Fünffach-Drehung her. Die drei Gewichte:
\begin{center}
\begin{tabular}{llll}
$p_0 = 2/9$ & $= 0{,}2222222$ & \textbf{rational} & --- kein $\varphi$ darin \\
$p_2 = (5 - 3\varphi)/9$ & $= 0{,}0162109$ & enthält $\varphi$ & \\
$p_1 = (2 + 3\varphi)/9$ & $= 0{,}7615669$ & enthält $\varphi$ & \\
\cmidrule{2-2}
Summe & $= 1{,}0000000$ & & \\
\end{tabular}
\end{center}

Bemerkenswert: der Anteil, der in die triviale Mode übergeht, ist
\textbf{rational}. $\varphi$ steckt allein in der Aufteilung der beiden
nicht-trivialen Moden. Alle drei tragen den Nenner $9 = 3^2$, dieselbe
Z$_3$-Struktur, in der $2/9$ als Orbit-Adresse sitzt.

\section{\texorpdfstring{$\varphi$}{phi} ist hier keine Wahl, sondern
ein Eigenwert}

Bei cut-and-project erscheint $\varphi$, weil der Beobachter die
Richtung so gelegt hat. Hier ist $\varphi$ der Eigenwert der Struktur
selbst:
\[
  \varphi^2 = \varphi + 1 .
\]
A$_5$ ist die Ikosaedergruppe. Ihre Fünffach-Achsen sind die zwölf
Ecken mit Koordinaten $(0, \pm1, \pm\varphi)$ und zyklisch. \textbf{Die
Richtungen sind durch $\varphi$ vollständig festgelegt} --- es bleibt
kein Winkel wählbar, nur welche der sechs Achsen, und A120 zeigt, dass
das eine diskrete Zweifachwahl ist ($2/9$ oder $4/9$), kein freier
Parameter.

Die Eindeutigkeit ist dort auch geprüft: 200 zufällige Drehachsen bei
$72^\circ$ geben Werte zwischen $0{,}036$ und $0{,}452$, keine einzige
trifft $2/9$. Bei dreizähliger Achse steht $2/5$, bei $144^\circ$ steht
$4/9$.

\section{Und das Fenster ebenfalls}

Ein Ikosaeder hat keinen Formparameter. Sein Umkugelradius folgt aus
der Kante:
\[
  |\text{Ecke}| = \sqrt{1 + \varphi^2} = 1{,}9021130, \qquad
  \frac{\text{Umkugel}}{\text{Kante}} = 0{,}9510565
  = \frac{\sqrt{\varphi\sqrt5}}{2}.
\]
Alle Größenverhältnisse sind durch $\varphi$ bestimmt. Was bleibt, ist
ein Maßstab --- und das ist genau die eine Größe, die A085 als den
Anker führt.

\section{Die Gegenüberstellung}

\begin{center}
\begin{tabular}{llll}
\toprule
 & Schnittrichtung & Fensterform & Fenstergröße \\
\midrule
cut-and-project, generisch & zwei freie Winkel & frei & frei \\
ikosaedrisch & 6 Achsen, diskret & das Ikosaeder, formfest & nur ein Maßstab \\
\bottomrule
\end{tabular}
\end{center}

Ein Kontinuum von Freiheiten auf der einen Seite, eine diskrete Auswahl
plus ein Maßstab auf der anderen.

\section{Was das bedeutet}

\begin{quote}
Wo $\varphi$ gewählt wird, ist es hineingelegt.\\
Wo $\varphi$ gefunden wird, legt es Richtung und Fenster fest.
\end{quote}

Das ist keine Wortspielerei, sondern der Unterschied zwischen einem
Parameter und einer Struktur. Eine Konstruktion, in der $\varphi$ als
Projektionsrichtung gesetzt wird, hat einen freien Winkel. Eine, in der
$\varphi$ als Eigenwert einer Gruppe auftritt, hat eine diskrete
Auswahl.

Und es kehrt den Prüfstein um. Statt nur zu fragen, \emph{wer die
Richtung gewählt hat}, lässt sich fragen:

\begin{quote}
\textbf{Gibt es eine Struktur, deren Eigenwerte die Richtung
erzwingen?}
\end{quote}

Kommt darauf eine Antwort, ist die Richtung keine Setzung mehr.

\chapter{Zwei Bedingungen statt drei Wege}

Wenn $\varphi$ Richtung und Fenster erzwingt, kostet Weg~2 keine freien
Parameter mehr --- so wenig wie Weg~1 mit gebundenem Radius. Die
Einteilung in drei Wege wird damit hinfällig, aber nicht ganz: eine
zweite Bedingung bleibt, und sie ist von $\varphi$ unabhängig.

\section{Was übertragbar wird}

\textbf{Die Freiheitsbilanz.} Erzwingt eine Struktur die Zuordnung, ist
sie keine Setzung mehr, gleich ob man sie Kompaktifizierung oder
Projektion nennt.
\begin{center}
\begin{tabular}{ll}
Weg 1 mit freiem Radius & ein Parameter \\
Weg 1 mit $\tilde T \cdot m = 1$ & keiner \\
Weg 2 mit freier Richtung & mehrere Parameter \\
Weg 2 mit ikosaedrischer Achse & keiner \\
\end{tabular}
\end{center}

Die Unterscheidung verläuft also nicht zwischen den Wegen, sondern
zwischen \emph{erzwungen} und \emph{gesetzt}.

\section{Was nicht übertragbar wird}

\textbf{Die Erhaltung der Information.} Eine Projektion ist per
Definition nicht injektiv. Auch bei erzwungener Richtung verschmelzen
$(1,0,0,+1)$ und $(1,0,0,-1)$ zum selben Bild $(1,0,0)$. Zwölf
Wicklungsvektoren werden zu sechs Bildern, gleichgültig wer die
Richtung festgelegt hat.

Beim Kompaktifizieren geschieht das nicht. Die Wicklungszahl bleibt
unterscheidbar --- $+1$ und $-1$ sind verschiedene Zustände und tragen
verschiedene Quantenzahlen, obwohl der Raum drei Dimensionen hat.

\section{Wo die beiden doch zusammenkommen}

Die Information verschwindet beim Kompaktifizieren nicht, sie
\textbf{wandert}. Vom Ort in einen Index. Was vorher eine vierte
Koordinate war, ist danach eine Wicklungszahl.

Eine Projektion mit erzwungener Richtung tut dasselbe --- \textbf{sofern
das Urbild-Paar einen erhaltenen Index trägt}. Trägt es keinen, ist es
Verlust. Trägt es einen, ist es Umbuchung, und dann ist der Unterschied
zur Kompaktifizierung nur noch einer der Sprechweise.

\section{Die Regel}

\begin{quote}
Nicht drei Wege, sondern zwei Bedingungen.

\textbf{1. Ist die Zuordnung durch eine Struktur erzwungen oder vom
Beobachter gesetzt?}

\textbf{2. Wird die verlorene Ortsinformation von einem erhaltenen
Index aufgefangen?}

Sind beide erfüllt, ist gleichgültig, wie man das Verfahren nennt.
Fehlt die erste, hat man Parameter. Fehlt die zweite, hat man Verlust.
\end{quote}

\section{Was das für die Beispiele heißt}

\begin{center}
\begin{tabular}{lll}
\toprule
 & Bedingung 1 & Bedingung 2 \\
\midrule
Kaluza-Klein, freier Radius & nein & ja \\
cut-and-project, generisch & nein & nein \\
FFGFT, $\tilde T \cdot m = 1$ & ja & ja \\
ikosaedrische Projektion & ja & nur mit Index \\
\bottomrule
\end{tabular}
\end{center}

Die letzte Zeile ist die interessante. Sie ist der Fall, in dem
$\varphi$ die Setzung beseitigt, aber die Frage nach dem Index
offenbleibt --- und die lässt sich nicht durch Geometrie beantworten,
sondern nur dadurch, dass man angibt, welche Größe den Index trägt.

In FFGFT ist das die Wicklungszahl auf der Massenrichtung. In einer
reinen Projektionskonstruktion müsste die entsprechende Größe benannt
werden, sonst ist der Verlust real.

\chapter{Der Befund, der die Sache entscheidet}

D4 muss gar nicht reduziert werden. Es zerfällt von selbst.

Von den 24 kürzesten Vektoren:
\begin{center}
\begin{tabular}{lcl}
\toprule
 & Anzahl & Bedeutung \\
\midrule
mit $x_4 = 0$ & \textbf{12} & liegen vollständig in den ersten drei Koordinaten \\
mit $x_4 \neq 0$ & \textbf{12} & sechs mit $x_4 = +1$, sechs mit $x_4 = -1$ \\
\bottomrule
\end{tabular}
\end{center}

Und die zwölf mit $x_4 = 0$ sind nicht irgendetwas:
\[
  \text{Kusszahl von D3} = \text{FCC} = 12 .
\]

FCC steht für \emph{face-centered cubic}, deutsch
kubisch-flächenzentriert oder kfz --- die dichteste Kugelpackung in
drei Dimensionen, Schichtenfolge ABCABC. Eine Kugel im Inneren berührt
sechs Nachbarn in ihrer Schicht, drei darüber und drei darunter:
$6 + 3 + 3 = 12$. Dass das Maximum ist, geht auf die Frage von Newton
und Gregory zurück und wurde 1953 endgültig bewiesen.

Gitter und Packung sind dasselbe Objekt in zwei Sprachen: die kürzesten
Vektoren von D3 sind $(\pm1, \pm1, 0)$ und Permutationen, genau zwölf,
und sie zeigen zu den zwölf Berührpunkten.

\section{Warum die Pyramide in drei Dimensionen grundlegend ist}

Die zwölf Nachbarvektoren bilden als konvexe Hülle ein
\textbf{Kuboktaeder} --- acht Dreiecke, sechs Quadrate, alle Ecken im
Abstand $\sqrt2$. Das ist die Nachbarschaftsfigur des dreidimensionalen
Raums.

Und dieselbe Packung hat eine zweite Gestalt, die jeder kennt. Legt man
Kugeln in die Mulden einer quadratischen Schicht --- der Orangenstapel
---, so hat jede Kugel im Inneren
\[
  4 \text{ Nachbarn unten} + 4 \text{ in der Ebene}
  + 4 \text{ oben} = 12
\]
Genau die Kusszahl. \textbf{Die quadratische Pyramide ist FCC}, nur
anders hingestellt: der Blick geht entlang der Würfeldiagonale statt
entlang der Kante. Hexagonale Schichten in der Folge ABCABC und
quadratische Schichten mit Mulden beschreiben dasselbe Gitter.

Das Fundamentale daran ist die Einzigkeit. In drei Dimensionen gibt es
genau \textbf{eine} dichteste Packung --- die Keplersche Vermutung, von
Hales 1998 gezeigt und 2014 formal verifiziert. Die Pyramide ist also
keine Bauform unter mehreren, sondern die Form, die die Dichte
erzwingt. Deshalb kristallisieren Kupfer, Gold und Aluminium so, und
deshalb funktioniert der Obststand.

Für diese Notiz zählt daran zweierlei. Erstens ist die 12 keine
Konventionszahl, sondern in 3D bewiesen maximal. Zweitens hat der
dreidimensionale Raum damit eine ausgezeichnete Nachbarschaftsfigur, an
der sich messen lässt, ob eine Konstruktion in ihm rechnet oder nicht:
wer räumlich mit 24 arbeitet, arbeitet nicht im Kuboktaeder.
\[
  24 = 12 + 12
\]

\textbf{Die Hälfte der D4-Kusszahl ist der dreidimensionale Raum. Die
andere Hälfte führt aus ihm heraus} --- und wird zu Wicklungsnachbarn,
sobald die vierte Richtung kompakt ist, je sechs in jede
Umlaufrichtung.

Das ist die Gitterform von
$\text{T}^4 = \text{T}^3 \times S^1_m$: drei Raumkreise plus die
vierte Richtung, und die Aufteilung $12 + 12$ fällt heraus statt
hineingelegt zu werden.

\chapter{Was daraus für die Kennzahlen folgt}

Die Kennzahlen von D4 sind \textbf{4D-Größen}. Sie sind nur dann
Raumgrößen, wenn alle vier Richtungen Raum sind.

\begin{center}
\begin{tabular}{lll}
\toprule
 & D4 (4D) & D3 $=$ FCC (3D) \\
\midrule
Kusszahl & 24 & 12 \\
Packungsdichte & $\pi^2/16 = 0{,}616850$ & $\pi/(3\sqrt2) = 0{,}740480$ \\
\bottomrule
\end{tabular}
\end{center}

Wer mit 24 als räumlicher Kusszahl rechnet, behauptet damit implizit
vier Raumdimensionen. Wer drei Raumdimensionen plus eine andere
Richtung hat, rechnet räumlich mit 12, und die anderen zwölf Nachbarn
sind etwas anderes --- keine Nachbarn im Raum, sondern Wicklungen.

Das ist keine Kleinigkeit. Es betrifft jede Größe, die aus der
Nachbarschaftszahl gebildet wird.

\section{Welche Packungsdichte gehört in \texorpdfstring{$K^{-36}$}{K hoch -36}?}

Daraus folgt eine Frage an die eigene Konstruktion, und sie musste
geprüft werden: der Bulk-Exponent der Sektorleiter (A270) ist an eine
Packungsdichte gebunden. Wenn nur drei der vier Richtungen Raum sind
--- ist es dann die 4D-Dichte oder die 3D-Dichte?

\begin{center}
\small
\begin{tabular}{lllll}
\toprule
 & Kehrwert der Dichte & verlangter Exponent
 & Abstand zu $k^*/100$ & zur Ganzzahl \\
\midrule
D4 (4D) & $16/\pi^2 = 1{,}6211389$ & \textbf{35{,}9926}
  & $0{,}098$ = \textbf{0{,}27\,\%} & $0{,}0074$ \\
D3 $=$ FCC (3D) & $3\sqrt2/\pi = 1{,}3504745$ & $22{,}3836$
  & $13{,}707$ = \textbf{38{,}08\,\%} & $0{,}3836$ \\
\bottomrule
\end{tabular}
\end{center}

Und bei ganzzahligem Exponent:
\begin{center}
\begin{tabular}{lllr}
D4: & $(75/74)^{36} = 1{,}621301$ & gegen $1{,}621139$ & $+0{,}010\,\%$ \\
D3: & $(75/74)^{22} = 1{,}343538$ & gegen $1{,}350474$ & $-0{,}514\,\%$ \\
\end{tabular}
\end{center}
(In der Schreibweise $K^{-n}$ mit $K = 74/75$; ausgeschrieben als
$(75/74)^n$, um das Vorzeichen des Exponenten eindeutig zu halten.)

\textbf{Das Ergebnis ist eindeutig.} Die 3D-Dichte passt nicht. Sie
verlangt Exponent $22{,}38$ statt $36{,}09$, also $38\,\%$ daneben, und
sie liegt zusätzlich nahe an einer Rundungsgrenze --- $0{,}38$ gegen
$0{,}007$ bei der 4D-Dichte. Selbst als isolierter Treffer wäre sie
schwach. \textbf{In $K^{-36}$ gehört die 4D-Dichte.}

\subsection*{Warum das kein Widerspruch ist}

Der Bulk ist T$^4$, nicht T$^3$. Die vierte Richtung trägt zur Packung
bei, auch wenn sie kein Raum ist --- weil dort gepackt wird, wo der
Torus ist, und der hat vier Richtungen.

Das ist zulässig, aber es muss gesagt werden. Die Massenrichtung ist
eine echte Torusrichtung mit ganzzahligen Wicklungen, keine
Beschriftung. Sie zählt bei der Packung mit, weil sie geometrisch
vorhanden ist, und sie zählt bei der räumlichen Kusszahl nicht mit,
weil sie kein Raum ist.

Die beiden Aussagen stehen also nebeneinander, ohne sich zu
widersprechen:
\begin{center}
\begin{tabular}{lcl}
räumliche Nachbarschaft & $\longrightarrow$ & 12 (FCC, die $x_4 = 0$ Vektoren) \\
Packung des Bulk & $\longrightarrow$ & $\pi^2/16$ (alle vier Richtungen) \\
\end{tabular}
\end{center}

\subsection*{Und was daraus für vergleichbare Konstruktionen folgt}

Dieselbe Zahl, zwei verschiedene Rechtfertigungen --- und
typischerweise ist nur eine davon ausgesprochen.

In FFGFT ist $\pi^2/16$ die Packungsdichte eines vierdimensionalen
Torus, von dem drei Richtungen Raum sind und eine die Masse trägt. Das
ist eine Aussage, die begründet werden kann und begründet wird.

In einer Konstruktion, die auf demselben Gitter rechnet, aber alle vier
Richtungen als räumlich führt, ist $\pi^2/16$ die Packungsdichte eines
vierdimensionalen \emph{Raums}. Das ist eine Behauptung über die Zahl
der Raumdimensionen --- und sie wird selten ausgesprochen, weil die
Zahl als neutrale Gitterkonstante erscheint.

Sie ist nicht neutral. Wer sie verwendet, sagt damit etwas über die
Zahl der Raumdimensionen, ob er will oder nicht.

\chapter{Anwendung: Konstruktionen mit vier räumlichen
Gitterrichtungen}

Der Fall, der in der Praxis vorkommt: eine Konstruktion rechnet auf D4
mit 24 Offsets, hat einen Schrittzähler für die Iteration und erklärt
nirgends eine Zeitrichtung. Daraus folgt zwingend: \textbf{alle vier
Koordinaten sind Raum.}

Das ist eine Behauptung über die Welt, nicht nur über das Modell ---
ein Objekt in vier Raumdimensionen ist kein Teilchen in unserem Raum,
sondern ein Gebilde in einem Raum, den es nicht gibt. Selbst bei
stimmender Knotenzahl bliebe offen, was sie mit einer Masse zu tun hat.

Und die Auswege sind unbequem:
\begin{itemize}
\item Eine Richtung wird nachträglich Zeit $\to$ dann bleiben
  \textbf{drei} Raumdimensionen, und die räumliche Kusszahl ist 12,
  nicht 24. Alles, was auf 24 aufsitzt, wäre neu zu rechnen.
\item Zeit wird als fünfte Richtung angehängt $\to$ dann sind es vier
  Raumdimensionen, also eine zu viel.
\item Es bleibt bei $4+0$ $\to$ dann ist zu sagen, wie ein
  vierdimensionales räumliches Objekt eine Masse in unserem Raum
  bekommt.
\end{itemize}

Das ist die grundlegende Frage an jede solche Konstruktion, und sie
geht den Zahlenbefunden voraus: solange sie offen ist, betrifft jede
Korrektur an einzelnen Größen nur Symptome.

\textbf{Sie wird meist nicht gestellt.} Wo die Zeit weder im Code noch
in den Dokumenten noch in der Korrespondenz vorkommt, existiert die
Frage für die Beteiligten noch nicht --- und das ist verständlich, denn
eine Iteration \emph{fühlt} sich wie Zeit an, ohne eine zu sein.

\textbf{Gestellt, nicht beantwortet.} Wer eine solche Konstruktion
baut, kennt ihre Interna; diese Notiz liest von außen. Möglicherweise
gibt es Antworten, die von hier aus nicht sichtbar sind.

\chapter{Die beiden Bedingungen, angelegt an die Stringtheorie}

Die Regel aus dem Kapitel \enquote{Zwei Bedingungen} lässt sich
anlegen, und das Ergebnis fällt nicht pauschal aus. Es benennt genau
eine Stelle.

\section{Bedingung 2 ist vorbildlich erfüllt}

Die kompakten Richtungen sind kompakt. Es wird nichts projiziert,
nichts verschmilzt. Wicklungszahlen, Kaluza-Klein-Impuls,
Windungssektoren, Hodge-Zahlen, Schnittzahlen --- alles erhaltene
Indizes. Die Ortsinformation der sechs eingerollten Richtungen wandert
vollständig in Quantenzahlen.

Das ist die Stärke des Ansatzes und sollte nicht übergangen werden. In
der Sprache der Regel: die Information wird umgebucht, nicht verloren.

\section{Bedingung 1 ist zur Hälfte erfüllt}

Die Trennlinie verläuft mitten durch die Konstruktion.

\textbf{Erzwungen:}
\begin{center}
\begin{tabular}{l}
$D = 10$ aus Anomaliefreiheit und kritischer Dimension \\
$N = 1$ SUSY in 4D verlangt SU(3)-Holonomie, also Calabi-Yau \\
daraus: $10 - 4 = 6$ kompakte Richtungen \\
\end{tabular}
\end{center}

Die \textbf{Zahl} der eingerollten Richtungen ist nicht wählbar. Das
ist eine echte Erzwingung, und sie ist stärker als alles, was FFGFT für
seine vier Kreise anführen kann --- dort ist die Wahl des T$^4$ als
[SETZUNG] markiert.

\textbf{Gesetzt:}
\begin{center}
\begin{tabular}{ll}
welche Calabi-Yau-Mannigfaltigkeit & ($\sim 10^5$ bekannte, Gesamtzahl unbekannt) \\
Flussquantenzahlen & ($\sim 10^{500}$ Kombinationen) \\
Kähler- und Strukturmoduli & (Kontinuum, dann stabilisiert) \\
Standardeinbettung oder nicht & \\
\end{tabular}
\end{center}

Die Zahl der Richtungen ist erzwungen. Die \textbf{Geometrie darin}
nicht.

\section{Was die Landschaft damit ist}

\begin{quote}
Die Landschaft ist genau der unerfüllte Teil von Bedingung 1.
\end{quote}

Nicht ein Fehler und nicht ein Skandal, sondern der Ort, an dem die
Struktur aufhört zu erzwingen. Eine präzise Ortsangabe statt eines
Vorwurfs.

Und sie erklärt, warum das Problem so hartnäckig ist: es ist nicht
durch mehr Rechnen zu beheben. Wo eine Struktur nicht erzwingt, kann
keine Rechnung sie dazu bringen. Es braucht eine zusätzliche Bedingung
von außen --- und genau danach wird seit Jahrzehnten gesucht.

\section{Die Asymmetrie, die beides verbindet}

Kompaktifiziert werden \textbf{Raum}richtungen. Die Zeit bleibt außen.
Nach A060 folgt daraus: kein dynamischer Maßstab im Zustandsraum. Alle
Skalen müssen von woanders kommen, und in dieser Konstruktion kommen
sie aus den Moduli.

Moduli sind aber genau das, was Bedingung 1 nicht erfüllt.
\begin{center}
\begin{tabular}{l}
Zeit außen \ $\to$ \ keine Skala in der Struktur \\
\phantom{Zeit außen} $\to$ \ Skalen aus Moduli \\
\phantom{Zeit außen} $\to$ \ Moduli sind der ungezwungene Teil \\
\phantom{Zeit außen} $\to$ \ Landschaft \\
\end{tabular}
\end{center}

Die Landschaft entsteht dort, wo die Skalen herkommen müssen. Und die
Skalen müssen von dort kommen, weil die Zeit draußen steht. Es ist eine
einzige Wahl, die beides erzeugt: sechs Raumrichtungen einrollen und
die Zeit offen lassen.

\section{Was das nicht sagt}

Es sagt nicht, dass die Stringtheorie falsch ist. Sie erfüllt
Bedingung~2 besser als alles hier Betrachtete, und ihre Dimensionszahl
ist zwingender begründet als die vier Kreise der FFGFT.

Es sagt auch nicht, dass die Alternative billiger wäre. Eine
eingerollte Zeit kostet Kausalitätseinwände, Paulis Theorem und die
thermische Lesart --- und sie ist bisher nirgends anerkannt.

Was es sagt, ist ausschließlich: \textbf{die Landschaft und die
Notwendigkeit angehefteter Skalen sind nicht zwei Probleme, sondern
eines.} Wer eines lösen will, muss beim anderen ansetzen.

\chapter{Prüfskript}

\begin{verbatim}
import itertools, math, collections

D4 = [v for v in itertools.product([-1,0,1], repeat=4)
      if sum(x*x for x in v) == 2]
D3 = [v for v in itertools.product([-1,0,1], repeat=3)
      if sum(x*x for x in v) == 2]

print(len(D4))                                 # 24
print(sum(1 for v in D4 if v[3] == 0))         # 12
print(sum(1 for v in D4 if v[3] != 0))         # 12
print(len(D3))                                 # 12  (FCC)
print(len(set(v[:3] for v in D4)))             # 18  (Projektion)
print(collections.Counter(
    v[3] for v in D4 if v[3] != 0))            # {1: 6, -1: 6}

print(math.pi**2 / 16)                         # 0,616850  D4
print(math.pi / (3 * 2**0.5))                  # 0,740480  D3
\end{verbatim}

Erweiterte Fassung mit allen Zahlen dieses Dokuments:
\texttt{python/\allowbreak Dok311\_\allowbreak Skripte/\allowbreak
ffgft\_\allowbreak 311\_\allowbreak vier\_\allowbreak auf\_\allowbreak
drei.py}.

\chapter{Bezüge}

\textbf{Interne Dokumente:} A060 (Zeit als Maßstab), A085 (Anker),
A120 ($\theta = 2/9$, Fünffach-Drehung), A270 (Bulk-Exponent 36, P35),
Dok.~275 (Skalenrekursionstiefe $k^*$), Dok.~310 (Resonanzgeometrie als
Referenz, Kammerton-Anker), Register 190 (R67: Unsicherheit des Faktors
100 im Nenner von $k^*/100$).

\medskip
\noindent\textbf{Zum Status:} Arbeitsnotiz für den eigenen Gebrauch.
Die Kapitel zu den Wegen, zum Zerfallsbefund $24 = 12+12$ und zur
Packungsdichte sind gerechnet; das Kapitel \enquote{Anwendung} ist eine
Frage und keine Feststellung.

\end{document}
